\documentclass[10pt]{article}
\usepackage[margin=20mm]{geometry}
\usepackage{amsmath,amssymb,mathtools,fontspec,unicode-math,xurl,hyperref}
\IfFontExistsTF{DejaVu Serif}{\setmainfont{DejaVu Serif}}{\setmainfont{Latin Modern Roman}}
\setmathfont{latinmodern-math.otf}
\setlength{\parindent}{0pt}\setlength{\parskip}{6pt}\sloppy
\title{Scientific initial review: M25-0003
Complete classification of global solutions
to the obstacle problem}
\author{}
\date{}
\begin{document}
\maketitle
Abstract: This review concerns the 45-page author manuscript arXiv:2208.03108v2, dated 19 August 2022, by Simon Eberle, Alessio Figalli and Georg S. Weiss. A false power-integral estimate enters the core paraboloid-envelope argument. A complete replacement estimate is supplied here; the classification theorem is not disproved. Two additional local statement defects concern an unnormalized uniqueness claim and the zero solution in a general blow-down lemma. The findings belong to the manifest-bound v2. A distinct author-hosted copy was checked selectively for version comparison; no equivalence to the journal Version of Record has been established.\par
\section*{Object, versions, and review coverage}
The frozen journal identity is Annals of Mathematics 201 (2025), 167–224, DOI 10.4007/annals.2025.201.1.3 [S2]. The publisher lists revision on 8 November 2023 and acceptance on 7 October 2024. The arXiv history inspected lists v1 and v2, with v2 dated 19 August 2022 [S1]. Bounded exact-title and identifier searches, including an authoritative-domain check for correction and erratum, identified no separate correction notice. This does not establish that no correction exists.\par
An author-hosted 45-page PDF [S3] differs from v2: it includes Appendix A and strengthens Definition 2.8(ii). Its HTTP Last-Modified field is 10 October 2024, which is transport metadata, not proof of its scientific revision date or journal-final status. The inspected passages retain the power-estimate chain and the unnormalized uniqueness wording; the choice of \(\delta\) on p. 14 is changed from equality to \(\delta\leq1/3\). These observations do not establish a correction of F01 or F02. The reviewed object remains v2.\par
All 45 v2 PDF pages, including references, were read. Active author TeX was inspected statically at the stated definitions, estimates and proof dependencies. The bound source filename is `Classification-of-global-solutions-in-the-obstacle-problem—V3.03.tex`; line references below refer to that file. Source pages 3, 6, 11–12, 14, 22, and 27–28 were also visually checked. No author code was executed. The input PDF, source file and manifest bindings were hash-verified.\par
The principal problem is \(\Delta u=\chi_{\{u>0\}}\), \(u\geq0\), on \(\mathbb R^N\). The monotone theorem requires all four conditions of Definition 2.8: nonempty contact interior, placement in a half-space with the origin on its boundary, \(\partial_Nu\leq0\), and a transverse positive-definite quadratic blow-down of trace one. These hypotheses are retained throughout the analysis. The distinction between \(N\geq4\) and \(N=3\) is essential.\par
\section*{F01: bounded and logarithmic power-integral regimes are lost}
Classification: Major; demonstrated invalid derivation on a core proof path.\par
\subsection*{Exact location}
Equations (3.5)–(3.6), pp. 11–12, estimate integrals over \[T_\delta=\{(y',t):t\geq0,\ |y'|^2<t^{1+\delta}\},\qquad 0<\delta<1.\] In (3.6) the integral of \(t^{(N-1)(1+\delta)/2+1-N}\) is bounded by a constant times \((1+|x|)^{(N-1)(1+\delta)/2+2-N}\), without a bounded-term or logarithmic case. See source lines 908–923. The proof of Proposition 4.1, p. 14, explicitly invokes (3.5) and (3.6) for its middle integral and then chooses \(\delta=1/3\) to infer \(V_{\mathcal C}(x)\geq-C(1+|x|)\); see lines 1058–1067. The lower-cutoff symbol \(B_R\) in that display is not defined there. The preceding decomposition uses the fixed ball \(B_{2\widehat r}\). The failure below already occurs in the unambiguous estimate (3.6), independently of this notation problem.\par
\subsection*{An explicit failure at the selected exponent}
Take \(N=4\) and \(\delta=1/3\). The exponent of the one-dimensional integrand in (3.6) is \(-1\), whereas the exponent of its proposed bound is zero. For fixed \(a>0\), \[\int_a^R\frac{dt}{t}=\log(R/a)\longrightarrow\infty.\] The same obstruction holds for the actual geometric integral, not only for an overly rough upper bound. Choose \(a>\max\{2\widehat r,1\}\) and \(R>a\), and restrict to \[a<t<R,\qquad |y'|<t^{2/3}.\] This region lies in \(T_{1/3}\cap(B_{2R}\setminus B_{2\widehat r})\), because \(t\leq|y|\leq\sqrt2\,t\) for \(t\geq1\). Its cross-sectional volume is \(|B_1^3|t^2\). Consequently, \[\int_{T_{1/3}\cap(B_{2R}\setminus B_{2\widehat r})}|y|^{-3}\,dy
 \geq \frac{|B_1^3|}{2^{3/2}}\log(R/a).\] No constant independent of \(R\) can satisfy (3.6) in this case. In the Proposition 4.1 majorant the moment is multiplied by \(|x|=R\), producing an \(R\log R\) term. It cannot be replaced by \(CR\) by the printed estimate.\par
There is also a bounded-integral error when the proposed power is negative. For example, \(N=4,\delta=1/4\) gives the integrand \(t^{-9/8}\). Its integral from a fixed positive lower limit tends to a strictly positive constant, whereas the printed power is \((1+R)^{-1/8}\to0\). Equation (3.5) has the corresponding issue in higher dimensions. These failures concern the explicit auxiliary integrals, rather than numerical rounding or a disputed coincidence set. No claim is made that \(T_{1/3}\) itself is a coincidence set.\par
\subsection*{A complete replacement estimate}
The defect is repairable without assuming Proposition 4.1 in advance. For its \(N\geq4\) proof choose \(\delta=1/4\) in Lemma 2.10 and set \(a_N=5(N-1)/8\). With a fixed lower cutoff at least one, the first-moment integral is bounded because \[a_N+1-N=-\frac{3(N-1)}8\leq-\frac98,\qquad
\int_a^{2R}t^{a_N+1-N}\,dt\leq C.\] The potential integral and tail obey \[\begin{aligned}
\int_a^{2R}t^{a_N+2-N}\,dt
 &\leq C(1+R^{7/8}),\\
R^2\int_R^\infty t^{a_N-N}\,dt
 &= C_N R^{a_N+3-N}\leq C_N R^{7/8}.\end{aligned}\] The first inequality follows by comparing the integrand to \(t^{-1/8}\) for \(t\geq1\), and remains valid when the actual integral is bounded. Retaining the bounded near contribution and multiplying the first moment by \(R\) gives \[V_{\mathcal C}(x)\geq-C(1+R)-CR^{7/8}\geq-C'(1+R),
\qquad R=|x|\geq1.\] The expansion \(u=p+V_{\mathcal C}\) and transverse coercivity \(p(x')\geq c_p|x'|^2\) then yield, for \(x\in\mathcal C\), \[c_p|x'|^2\leq C(1+|x'|+x_N).\] Young’s inequality absorbs the linear \(|x'|\) term and gives \(|x'|^2\leq C'(1+x_N)\), which is the needed envelope.\par
The earlier subquadratic estimate in Lemma 3.3 can also be justified without the false intermediate powers. For \(N\geq3\) and \(t\geq1\), compare the three relevant integrands directly to the \(N=3\) powers \(t^\delta,t^{\delta-1},t^{\delta-2}\) before integrating. The first two contributions are \(O(1+R^{1+\delta})\), and the tail is \(O(R^{1+\delta})\). The translated singular kernel integrated over transverse balls contributes \(O(R^{(3+\delta)/2})\). Setting \(\delta=1/2\) gives \(O(1+R^{7/4})\). This validates the regularized-potential construction and its subsequent Liouville expansion at this step.\par
\subsection*{Impact and counterevidence}
Proposition 4.1 supplies the envelope needed by Lemma 5.1, the linear matching in Corollary 6.3, the ACF ordering in Proposition 7.1, and the \(N\geq4\) part of Theorem 1.2. The false estimate is therefore on a core derivation, warranting Major under the prescribed standard. The replacement above is strong counterevidence to any claim that the geometric theorem is false. It establishes a repair of this particular proof step; it does not certify all remaining dependencies. The separate \(N=3\) envelope proof is not refuted by this calculation.\par
\section*{F02: quadratic blow-down does not determine the scale of the ellipsoid}
Classification: Minor; a false uniqueness assertion not needed by the core construction.\par
Lemma 6.1, pp. 27–28, asserts uniqueness of an ellipsoid \(E'\) for the specified quadratic blow-down \(p\), without a volume, semiaxis, or other scale normalization. The same claim appears in Step 1 of its proof, source lines 1673–1709.\par
If \(P=P_{E'}\) is one such paraboloid and \(u_P\) its solution, the rescaling explicitly given in the lemma, \[u_{\gamma P}(x)=\gamma^2u_P(x/\gamma),\] satisfies the same obstacle equation and has the same quadratic blow-down \(p\), since \(p\) is homogeneous of degree two. But \(\gamma P_{E'}=P_{\sqrt\gamma E'}\); taking \(\gamma\ne1\) changes the ellipsoid. Thus the existence and scaling parts of the lemma are incompatible with its unnormalized uniqueness claim.\par
There is also a direct potential check. In dimension \(d=N-1\) let \(E'_a\) be the ball of radius \(a>0\). Its radial Newtonian potential \(V_a\) satisfies \(-\Delta V_a=1\) inside this ball. Regularity at the centre and the radial equation give \[V_a(r)-V_a(0)=-\frac{r^2}{2d},\qquad 0\leq r\leq a.\] This holds for every \(a\), including the logarithmic kernel for \(d=2\). Hence the identity cited in Step 1 cannot determine \(a\).\par
For an independent realization by actual three-dimensional obstacle solutions, fix \(s>0\), write \(\rho^2=x_1^2+x_2^2\), and define \[Z=x_3-s/2,\quad
q=\sqrt{\rho^2+Z^2},\quad
\xi=q-Z,\quad \eta=q+Z.\] For \(\xi>s\) put \[\begin{aligned}
g(\xi)&=\frac{s}{4}\{\xi\log(\xi/s)-\xi+s\},\\
h(\xi)&=\frac14\{\xi-s-s\log(\xi/s)\},&
u_s&=g(\xi)+\eta h(\xi),\end{aligned}\] and set \(u_s=0\) for \(\xi\leq s\). Both \(g\) and \(h\) are nonnegative for \(\xi\geq s\), and \(g(s)=h(s)=g'(s)=h'(s)=0\). The axisymmetric Laplacian gives \[\Delta u_s=
\frac4{\xi+\eta}
 \left\{(\xi g')'+h+\eta(\xi h')'\right\}=1
 \quad\text{where }\xi>s,\] because the expression in braces is \((\xi+\eta)/4\). The zero extension has no first-derivative jump and is locally \(C^{1,1}\). Its contact set is \(\rho^2\leq2s x_3\), while its quadratic blow-down is \(\rho^2/4\), independent of \(s\): \(\xi\eta=\rho^2\), and the remaining terms grow at most \(O(R\log R)\) on \(B_R\); inside the contact set \(\rho^2=O(R)\). Compactness upgrades this blow-down to the regularity used in the lemma. Thus \(s=1\) and \(s=2\) give different cross-sectional disks with the same \(p\).\par
Exact coefficient checks of the PDE identity were executed; Cartesian finite-difference checks were used only as a sanity check. The identity and boundary conditions supply the proof. The subsequent construction requires a choice of one \(E'\) and then deliberately varies \(\gamma\). It does not require absolute uniqueness. Delete that claim or specify a scale normalization and its appropriate uniqueness theorem. No breakdown of the classification argument from this extra claim is demonstrated.\par
\section*{F03: the zero solution is omitted from the general blow-down lemma}
Classification: Minor; an excluded degenerate case needs an explicit convention.\par
Lemma 2.3, p. 5, source lines 499 onward, is stated for a global nonnegative solution of (1.1), and lists only a half-space blow-down or a quadratic form with trace-one Hessian. The smooth global function \(u\equiv0\) satisfies \(\Delta u=0=\chi_{\{u>0\}}\) and \(u\geq0\), but all its quadratic rescalings are zero. Zero is neither listed alternative. The lemma needs \(u\not\equiv0\) or an additional zero alternative.\par
This case cannot enter the main monotone proof: Definition 2.8(ii) excludes its contact set \(\mathbb R^N\), and condition (iv) also excludes a zero blow-down. Theorem 1.1 should likewise make its nontrivial-solution convention explicit, or explicitly include the full space as a degenerate contact set. Whether a zero-dimensional base is permitted in the term cylinder is a convention; the counterexample to Lemma 2.3 is independent of it. This is a scope correction, not evidence of a new nontrivial geometric class missing from the result.\par
\section*{Checks of the remaining argument}
The regularized kernel was checked for sign, normalization, homogeneity and cancellation. Its value and first derivative vanish at \(x=0\), its far-field remainder is \(O(|x|^2|y|^{-N})\), and its distributional Laplacian is the negative source density. The harmonic difference \(u-p-V_{\mathcal C}\) is subquadratic, hence affine; the value and gradient at the origin eliminate that affine part. The first moment for a true paraboloid is integrable starting at \(N=4\), whereas the unregularized potential needs \(N\geq6\). The cross-sectional power count confirms these thresholds.\par
For the three-dimensional envelope, the chain from convexity of sections to concavity of \(\sqrt{H}\), the one-dimensionalized potential, the running maximum \(\Psi\), and differential comparison was checked. With \[h(z)=-Cz-z\log z,\qquad
G(t)=A\log t-A(C+\log A),\] one has \(h(G')=G/t\) exactly. For \(A=t_0^{2\delta}\) the coefficient of \(\log t_0\) in \(G(t_0)/t_0^{2\delta}\) is \(1-2\delta\), not the \(2-2\delta\) printed on p. 22. The required domination still holds for \(0<\delta<1/2\) and sufficiently large \(t_0\). Asymptotic thresholds can be enlarged after fixing \(A\). This local arithmetic slip does not refute the section-growth conclusion.\par
For differences of two obstacle solutions, \(w\Delta w\geq0\) yields subharmonic positive and negative parts and the Caccioppoli estimate. The ACF scaling was checked explicitly: \[\Phi\bigl(v(r\,\cdot)/r,\rho\bigr)=\Phi(v,r\rho).\] The weighted energy on \(B_\delta\) contributes the factor \(\delta^2\), as in (8.9), preventing concentration at the origin in the ACF limit. The strong-convergence proof needs only \(w_k\Delta w_k\geq0\) and a harmonic weak limit, so it applies to the linear rescalings even though their summands are no longer normalized obstacle solutions.\par
The zero-scale limit in Section 8 must be interpreted as the solution \(p\) specified by the text, rather than as the literal set product \(0P\). Its zero set is the vertical axis. That axis, as well as a half-axis, has locally zero Sobolev \(2\)-capacity in \(\mathbb R^3\): a transverse logarithmic cutoff has energy \(O(1/|\log\epsilon|)\). Using the axis in the removable-set argument therefore does not obstruct the harmonic-limit step. The two ACF alternatives, horizontal translation by \(Q^{-1}b'\), the \(R/\log R\) adjustment, and the ordering/continuity argument were traced to the final classification. In the zero-scale case the sign of \(b_3\) can be tested at \(-e_3/2\), inside the ball of convergence in (8.11).\par
For Section 9 the directional ACF argument forces a single sign in each null direction. A two-dimensional kernel would contain a direction orthogonal to the gradient at a fixed positive point, contradicting strict monotonicity. After choosing the orientation to agree with Definition 2.8, the remaining null direction gives the monotone case. These checks distinguish the intended proof from local sign or radius slips; no further nontrivial counterexample has been established.\par
\section*{Assessment and limits}
The findings are one Major derivation defect and two Minor statement defects for v2. A concrete repair is proved for the Major estimate. The other defects concern information not needed in the nontrivial monotone construction. Their combination does not demonstrate failure of the classification theorem.\par
All manuscript pages were read, and the algebraic, scaling, boundary, and dependency checks above were performed independently. The full external proofs of free-boundary regularity, the ACF equality case, the compact-coincidence classification, and the cited stability of quadratic blow-down under translated limits were not all rederived. The author-hosted comparison was selective, not a second full review. No journal-final comparison or whole-paper correctness certification is claimed. Coverage, findings, provenance, independent code and results, actual compiler/dependency hashes, logs, and every-page report inspection are supplied separately.\par
\section*{Compact conclusion table}
Item | Judgment | Repair and remaining decision\par
F01, potential estimate | Major; power-integral bound invalid | The report supplies a corrected fixed-\(\delta\) bound and lower-potential argument for the required paraboloid containment. No counterexample to the classification theorem.\par
F02, Lemma 6.1 | Minor; uniqueness fails under scaling | Retain the existence statement actually used. The explicit scaled family refutes uniqueness, not the subsequent existence construction.\par
F03, Lemma 2.3 | Minor; zero-solution omission | Include the zero solution among the general alternatives. Definition 2.8 excludes it from the later nontrivial monotone case.\par
Scope and decision | v2 initial review complete | Later author-copy comparison was selective; full journal-version equivalence remains unestablished. Repairs are proposed, not author-adopted. Editorial decision remains pending.\par
S3 S. Eberle, A. Figalli and G. S. Weiss, Complete classification of global solutions to the obstacle problem, arXiv:2208.03108v2 (19 August 2022), reviewed PDF and bound TeX. https://arxiv.org/abs/2208.03108v2. Annals of Mathematics, article metadata, 201 (2025), 167–224, DOI 10.4007/annals.2025.201.1.3. https://annals.math.princeton.edu/2025/201-1/p03. A. Figalli, author-hosted distinct manuscript copy; selective version-comparison evidence, retrieved 25 October 2026. https://people.math.ethz.ch/\textasciitilde{}afigalli/papers-pdf/Complete-classification-of-global-solutions-to-the-obstacle-problem.pdf.\par
\end{document}
